% CHƯƠNG 5 - Rubric: Phân tích, hiểu code thực nghiệm (2,5 điểm)
\chapter{Tái hiện thực nghiệm}
\label{chuong:taihien}

\section{Cấu trúc mã nguồn của tác giả}

\subsection{Mackey-Glass}

Mã nguồn cho thực nghiệm Mackey-Glass gồm bốn tệp chính. Tệp \texttt{utils/generate.py} sinh
(hoặc đọc từ \texttt{data.pkl} có sẵn) chuỗi Mackey-Glass thô. Tệp \texttt{utils/utils.py}
định nghĩa lớp \texttt{MackeyGlass} (sinh dữ liệu), hàm \texttt{create\_time\_series\_dataset}
(cắt chuỗi thành cặp đầu vào/nhãn theo lookback window, tầm dự báo và offset như
Bảng~\ref{bang:mg-offset}), lớp mô hình \texttt{RNN}, và hàm \texttt{KL} tính hàm mất mát chưng
cất. Tệp \texttt{exp/base\_exp.py} là điểm vào chính, thực hiện tuần tự ba giai đoạn huấn
luyện (Hình~\ref{hinh:quytrinh-fgl}) và in ra MSE trên tập kiểm tra của cả ba mô hình. Tệp
\texttt{exp/drift\_exp.py} là một thực nghiệm mở rộng của chính tác giả, kết hợp FGL với phát
hiện trôi khái niệm bằng kiểm định Page-Hinkley trong lúc suy luận.

\subsection{CHB-MIT}

Mã nguồn CHB-MIT phức tạp hơn, gồm các nhóm tệp: \texttt{utils/load\_signals\_\{teacher,
student\}.py} đọc tệp \texttt{.edf} bằng thư viện \texttt{mne} và cắt cửa sổ 30 giây;
\texttt{utils/prep\_data\_\{teacher,student\}.py} áp dụng biến đổi STFT, oversampling lớp
ictal, rồi lưu cache kết quả bằng \texttt{hickle} vào \texttt{cachedir-teacher/} hoặc
\texttt{cachedir-student/}; \texttt{models/model.py} định nghĩa hai kiến trúc
\texttt{CNN\_LSTM\_Model} và \texttt{MViT}; \texttt{exp/seizure\_detection.py},
\texttt{exp/seizure\_prediction.py} và \texttt{exp/FGL\_CHBMIT.py} là ba điểm vào tương ứng
với ba giai đoạn huấn luyện, lưu checkpoint vào \texttt{pytorch\_models/}.

\section{Kết quả tái hiện trên Mackey-Glass}

\subsection{Lần chạy chuẩn}

Với tham số mẫu (\texttt{horizon=5, num\_bins=50, alpha=0.5, epochs=20}, seed mặc định 42),
dùng trực tiếp \texttt{data.pkl} tác giả cung cấp, kết quả thu được trình bày trong
Bảng~\ref{bang:mg-chuan}.

\begin{table}[H]{Kết quả tái hiện lần chạy chuẩn trên Mackey-Glass (seed 42)}
	\label{bang:mg-chuan}
	\centering
	\begin{tabular}{lc}
		\toprule
		\textbf{Mô hình} & \textbf{MSE trên tập kiểm tra} \\
		\midrule
		Teacher (dự báo 1 bước) & 5,4391 \\
		Baseline (dự báo 5 bước, không FGL) & 139,3832 \\
		Student (dự báo 5 bước, có FGL) & 104,3942 \\
		\bottomrule
	\end{tabular}
\end{table}

Mức giảm MSE từ Baseline sang Student là $(139,3832 - 104,3942)/139,3832 \approx 25,1\%$, khớp
khá sát với mức \textbf{23,4\%} tác giả công bố (có loại outlier), ngay từ lần chạy đầu tiên
với đúng tham số mẫu, chưa chỉnh gì thêm.

\subsection{Sweep sáu seed để kiểm tra độ ổn định}
\label{sec:mg-sweep-seed}

Vì tác giả hardcode seed cố định (42) và không cung cấp cách đổi seed qua dòng lệnh, báo cáo
đã bổ sung tham số \texttt{--seed} để chạy lại trên sáu seed khác nhau, kết quả trong
Bảng~\ref{bang:mg-sweep}.

\begin{table}[H]{Kết quả sweep sáu seed trên Mackey-Glass}
	\label{bang:mg-sweep}
	\centering
	\small
	\begin{tabular}{ccccc}
		\toprule
		\textbf{Seed} & \textbf{Teacher} & \textbf{Baseline} & \textbf{Student} & \textbf{Giảm MSE} \\
		\midrule
		0  & 4,8331 & 156,1701 & 102,7514 & 34,20\% \\
		1  & 5,4642 & 136,5673 & 106,0795 & 22,32\% \\
		2  & 5,7293 & 143,5308 & 104,4262 & 27,25\% \\
		3  & 6,2236 & 140,3592 & 107,6203 & 23,32\% \\
		4  & 7,4506 & 147,5238 & 108,6868 & 26,33\% \\
		42 & 5,4391 & 139,3832 & 104,3942 & 25,10\% \\
		\midrule
		\textbf{TB} & 5,857 & 143,922 & 105,660 & \textbf{26,42\%} \\
		\bottomrule
	\end{tabular}
\end{table}

Toàn bộ sáu seed đều cho Student MSE thấp hơn Baseline MSE - FGL cải thiện nhất quán, không
có seed nào đảo chiều. Mức cải thiện trung bình 26,42\% (độ lệch chuẩn 3,86) \textbf{cao hơn
một chút} so với 23,4\% công bố, có thể do khác biệt nhỏ về cách loại outlier mà tác giả áp dụng mà thực nghiệm
này chưa làm, hoặc do bài báo lấy trung bình trên nhiều cấu hình thực nghiệm hơn. Dù vậy, kết
quả tái hiện khớp \textbf{đúng hướng và đúng độ lớn} với số liệu công bố.

\subsection{Thực nghiệm mở rộng của tác giả: quét tầm dự báo với Page-Hinkley}

Kịch bản \texttt{drift\_exp.py} (của tác giả, không phải tự đề xuất) quét tầm dự báo từ 2 đến
15 bước, với $\alpha = 0,0$ (student học hoàn toàn từ teacher, không dùng nhãn thật). Kết quả
trình bày trong Bảng~\ref{bang:mg-drift}.

\begin{table}[H]{Kết quả quét tầm dự báo 2-15 với Page-Hinkley trên Mackey-Glass}
	\label{bang:mg-drift}
	\centering
	\small
	\begin{tabular}{ccccc}
		\toprule
		\textbf{H} & \textbf{Teacher} & \textbf{Baseline} & \textbf{Student} & \textbf{Giảm MSE} \\
		\midrule
		2  & 0,4479 & 0,6943 & 0,7443 & $-7,20\%$ \\
		3  & 0,4458 & 1,5198 & 1,2870 & 15,32\% \\
		4  & 0,4542 & 2,5526 & 2,0479 & 19,77\% \\
		5  & 0,4516 & 3,5365 & 3,0359 & 14,15\% \\
		6  & 0,4453 & 4,3719 & 3,9307 & 10,09\% \\
		7  & 0,4573 & 4,9760 & 4,7771 & 4,00\% \\
		8  & 0,4562 & 5,2880 & 5,0052 & 5,35\% \\
		9  & 0,4589 & 5,5932 & 5,1766 & 7,45\% \\
		10 & 0,4469 & 4,9651 & 4,5422 & 8,52\% \\
		11 & 0,4365 & 4,0057 & 3,3552 & 16,24\% \\
		12 & 0,4380 & 3,1568 & 2,6109 & 17,29\% \\
		13 & 0,4359 & 2,7406 & 2,3344 & 14,82\% \\
		14 & 0,4401 & 2,3135 & 2,0281 & 12,34\% \\
		15 & 0,4469 & 2,1323 & 1,8453 & 13,46\% \\
		\midrule
		\textbf{TB} & 0,446 & - & - & \textbf{10,83\%} \\
		\textbf{TB (bỏ H=2)} & - & - & - & \textbf{12,22\%} \\
		\bottomrule
	\end{tabular}
\end{table}

13/14 giá trị horizon cho FGL cải thiện; chỉ \textbf{horizon=2 là ngoại lệ âm} ($-7,20\%$), tức
ở tầm dự báo rất ngắn, FGL kết hợp Page-Hinkley \textbf{làm tệ hơn} so với không dùng FGL.
Teacher MSE gần như hằng số ($\approx 0,44$-$0,46$) ở mọi horizon, hợp lý vì teacher luôn giải
bài toán 1 bước cố định. Baseline MSE không tăng đơn điệu theo horizon (tăng từ h=2 đến h=9,
rồi giảm từ h=10 đến h=15) - quan sát thú vị, có thể liên quan tới chu kỳ trễ đặc trưng
$\tau=17$ của hệ Mackey-Glass, không phải lỗi mã nguồn.

\section{Kết quả tái hiện trên CHB-MIT}
\label{sec:chbmit-ketqua}

Trên máy cá nhân, Giai đoạn 1 (Teacher, bệnh nhân 1) cho AUC-ROC $=0,28$ - thấp hơn một bộ
phân loại ngẫu nhiên (0,5), một kết quả bất thường; Giai đoạn 2 và 3 không hoàn thành do hết bộ
nhớ khi xây dựng cache đặc trưng STFT. Thực nghiệm sau đó được chạy lại trên GPU Kaggle, từ bản
sạch của repo gốc tác giả, mở rộng dần từ ba lên sáu bệnh nhân (1, 3, 5, 9, 13, 23) - chọn đa
dạng theo kết quả Bảng 1 của bài báo gốc: bệnh nhân 13 đạt AUC cao nhất (1,00), bệnh nhân 9 thấp
nhất (0,57), bệnh nhân 23 là trường hợp FGL duy nhất thua CNN-LSTM baseline trong bài báo. Tại
thời điểm viết báo cáo, sáu bệnh nhân hoàn thành các giai đoạn ở mức độ khác nhau (Bảng~\ref{bang:chbmit-tiendo}).

\begin{table}[H]{Tiến độ ba giai đoạn huấn luyện theo bệnh nhân}
	\label{bang:chbmit-tiendo}
	\centering
	\begin{tabular}{cccc}
		\toprule
		\textbf{Bệnh nhân} & \textbf{Teacher} & \textbf{Baseline} & \textbf{Student (FGL)} \\
		\midrule
		1  & Xong & Xong & Xong \\
		3  & Xong & Xong & Xong \\
		5  & Xong & Xong & Xong \\
		9  & Chưa & Chưa & Chưa \\
		13 & Xong & Chưa & Xong \\
		23 & Xong & Chưa & Chưa \\
		\bottomrule
	\end{tabular}
\end{table}

Bệnh nhân 9 chưa huấn luyện được ở bất kỳ giai đoạn nào dù đã áp dụng cơ chế tự phát hiện và xoá
cache hỏng (do hết dung lượng đĩa ở phiên Kaggle trước) rồi thử lại - nguyên nhân cụ thể còn
cần điều tra thêm, có thể không chỉ do hết đĩa. Bệnh nhân 13 và 23 mới có Teacher, Giai đoạn 2
(Baseline) của hai bệnh nhân này chưa chạy nên chưa so sánh được mức tăng của FGL.

\begin{table}[H]{AUC-ROC của Teacher (Giai đoạn 1) theo bệnh nhân, chạy trên Kaggle}
	\label{bang:chbmit-detection}
	\centering
	\begin{tabular}{cc}
		\toprule
		\textbf{Bệnh nhân} & \textbf{AUC-ROC} \\
		\midrule
		1 (trung bình 3 lần chạy) & 0,6451 \\
		3 & 0,9067 \\
		5 & 0,9911 \\
		13 & 0,7286 \\
		23 & 0,9822 \\
		\bottomrule
	\end{tabular}
\end{table}

Tất cả đều cao hơn rõ rệt mức ngẫu nhiên - khác hẳn kết quả bất thường (0,28) trên máy cá
nhân, cho thấy đó nhiều khả năng là một lần khởi tạo ngẫu nhiên không may
(\texttt{seizure\_detection.py} không cố định seed) hơn là lỗi hệ thống; nghi vấn \texttt{shuffle=False}
(mục~\ref{item:shuffle-false}) chưa bị bác bỏ hoàn toàn nhưng không còn đủ cơ sở khẳng định là
nguyên nhân chính.

\subsection{Baseline so với Student/FGL: hai lần chạy độc lập}
\label{sec:chbmit-hailanchay}

Cả \texttt{seizure\_prediction.py} và \texttt{FGL\_CHBMIT.py} không cố định seed ngẫu nhiên.
Trong quá trình mở rộng thực nghiệm, ba bệnh nhân 1, 3, 5 vô tình được huấn luyện lại hoàn toàn
độc lập ở một phiên Kaggle sau - cùng mã nguồn, cùng siêu tham số, chỉ khác lần khởi tạo
trọng số và việc chia dữ liệu ngẫu nhiên giữa các trial. Thay vì chỉ giữ lại lần chạy mới nhất,
báo cáo trình bày \textit{cả hai lần chạy} vì bản thân sự khác biệt giữa chúng là một bằng
chứng thực nghiệm có giá trị cho hạn chế đã nêu ở mục~\ref{sec:mg-sweep-seed}: kết quả CHB-MIT
cũng nhạy với lần khởi tạo ngẫu nhiên, không chỉ Mackey-Glass.

\begin{table}[H]{Lần chạy 1 - AUC-ROC trung bình Baseline và Student/FGL theo bệnh nhân}
	\label{bang:chbmit-sosanh-benhnhan}
	\centering
	\begin{tabular}{cccc}
		\toprule
		\textbf{Bệnh nhân} & \textbf{Baseline} & \textbf{Student (FGL)} & \textbf{Mức tăng} \\
		\midrule
		1 & 0,9173 & 0,9068 & $-1,1\%$ \\
		3 & 0,4276 & 0,7253 & $+69,6\%$ \\
		5 & 0,1963 & 0,3901 & $+98,7\%$ \\
		\midrule
		\textbf{Trung bình 3 bệnh nhân} & \textbf{0,5137} & \textbf{0,6741} & \textbf{+31,2\%} \\
		\bottomrule
	\end{tabular}
\end{table}

\begin{table}[H]{Lần chạy 2 (độc lập, cùng bệnh nhân) - AUC-ROC trung bình Baseline và Student/FGL}
	\label{bang:chbmit-sosanh-benhnhan-l2}
	\centering
	\begin{tabular}{cccc}
		\toprule
		\textbf{Bệnh nhân} & \textbf{Baseline} & \textbf{Student (FGL)} & \textbf{Mức tăng} \\
		\midrule
		1 & 0,6825 & 0,7071 & $+3,6\%$ \\
		3 & 0,3905 & 0,5667 & $+45,1\%$ \\
		5 & 0,0431 & 0,5695 & $+1221\%$\footnotemark \\
		\midrule
		\textbf{Trung bình 3 bệnh nhân} & \textbf{0,3720} & \textbf{0,6144} & \textbf{+65,2\%} \\
		\bottomrule
	\end{tabular}
\end{table}
\footnotetext{Baseline bệnh nhân 5 ở lần chạy 2 gần như suy biến: 2/3 trial có
FPR = Sensitivity = 0 (mô hình dự báo toàn bộ về một lớp), kéo AUC trung bình xuống 0,0431 -
gần mức một mô hình không học được gì. Với mẫu số gần 0, mức tăng phần trăm (+1221\%) không
còn nhiều ý nghĩa thống kê; nên đọc cùng với \textbf{giá trị AUC tuyệt đối} (0,57) chứ không chỉ
tỉ lệ tăng.}

Mỗi bệnh nhân chạy ba trial; kết quả giữa hai lần chạy khác biệt rõ rệt dù cùng mã nguồn và
siêu tham số: mức tăng trung bình trên ba bệnh nhân là \textbf{+31,2\%} ở lần 1 nhưng
\textbf{+65,2\%} ở lần 2 - chênh lệch hơn gấp đôi, chỉ do khác seed ngẫu nhiên. Bệnh nhân 1 đặc
biệt nhạy: Baseline dao động từ 0,92 (lần 1) xuống 0,68 (lần 2), và mức tăng của FGL đổi từ âm
nhẹ ($-1,1\%$) sang dương nhẹ ($+3,6\%$). Bệnh nhân 5 cũng cho thấy Baseline có thể suy biến gần
hoàn toàn ở một lần chạy (0,04) nhưng không phải lần khác (0,20) - CNN-LSTM baseline cho bài
toán dự báo (preictal vs.\ interictal, khó hơn phát hiện) dường như kém ổn định khi huấn luyện
với dữ liệu ít và không cố định seed.

So với mức 44,8\% bài báo công bố, cả hai lần chạy ba bệnh nhân đều \textbf{thấp hơn}, nhưng
theo hai hướng khác nhau (31,2\% và 65,2\%) - bản thân khoảng cách giữa hai con số này
(\textasciitilde34 điểm phần trăm, trên cùng ba bệnh nhân) đã lớn hơn khoảng cách giữa mỗi con
số với 44,8\%. Điều này củng cố thêm nhận định ở mục~\ref{sec:mg-sweep-seed} rằng, với số lần
lặp lại còn ít, kết quả tái hiện nên được đọc là ``khớp hướng'' (FGL cải thiện so với Baseline)
chứ chưa phải một ước lượng đáng tin cậy cho \textit{độ lớn} chính xác của mức cải thiện.

Bệnh nhân 13 - bệnh nhân có AUC cao nhất trong bài báo gốc (1,00) - cho AUC Student/FGL
trung bình \textbf{0,9798} (ba trial: 0,957; 0,983; 1,000), nhất quán cao và ổn định hơn hẳn so
với các bệnh nhân khác, dù chưa có Baseline để tính mức tăng tương ứng. Đây là quan sát cùng
hướng với bài báo gốc (bệnh nhân 13 dễ nhất), nhưng cần Giai đoạn 2 hoàn tất mới kết luận được
về mức cải thiện.

\subsection{Đối chiếu trực tiếp với số liệu công bố theo từng bệnh nhân}

Bài báo gốc công bố AUC-ROC của CNN-LSTM và FGL cho từng bệnh nhân (Bảng 1 trong bài báo), cho
phép đối chiếu trực tiếp thay vì chỉ so sánh mức trung bình. Bảng~\ref{bang:chbmit-doichieu-baibao}
chỉ liệt kê ba bệnh nhân có đủ cả Baseline và FGL để so sánh (1, 3, 5); bệnh nhân 9, 13, 23
chưa đủ dữ liệu (mục~\ref{sec:chbmit-ketqua}) nên không đưa vào.

\begin{table}[H]{Đối chiếu AUC-ROC tái hiện với số liệu công bố (bệnh nhân 1, 3, 5)}
	\label{bang:chbmit-doichieu-baibao}
	\centering
	\small
	\begin{tabular}{ccccc}
		\toprule
		& \multicolumn{2}{c}{\textbf{Bài báo}} & \multicolumn{2}{c}{\textbf{Tái hiện (L1 / L2)}} \\
		\textbf{Bệnh nhân} & \textbf{CNN-LSTM} & \textbf{FGL} & \textbf{Baseline} & \textbf{FGL} \\
		\midrule
		1 & 0,92 & 0,98 & 0,92 / 0,68 & 0,91 / 0,71 \\
		3 & 0,27 & 0,97 & 0,43 / 0,39 & 0,73 / 0,57 \\
		5 & 0,12 & 0,96 & 0,20 / 0,04 & 0,39 / 0,57 \\
		\bottomrule
	\end{tabular}
\end{table}

Hai quan sát nổi bật. \textbf{Thứ nhất}, Baseline tái hiện khớp khá tốt với CNN-LSTM công bố ở
bệnh nhân 1 (0,92 ở lần 1, trùng khít) và tương đối gần ở bệnh nhân 5 (0,20 so với 0,12 ở lần
1), trong khi FGL tái hiện lại lệch xa so với FGL công bố ở hầu hết các trường hợp (bệnh nhân 3:
0,57-0,73 so với 0,97; bệnh nhân 5: 0,39-0,57 so với 0,96) - ngoại lệ là bệnh nhân 13 ở trên,
nơi FGL tái hiện gần như trùng khít bài báo. \textbf{Thứ hai}, một khác biệt cấu hình cụ thể
được bài báo nêu rõ trong phần Methods có thể giải thích phần lớn độ lệch này: tác giả chọn
$\alpha$ \textit{tối ưu riêng cho từng bệnh nhân} qua một hyperparameter sweep, trong khi thực
nghiệm tái hiện này dùng $\alpha=0,5$ \textbf{cố định cho mọi bệnh nhân}. Ở bệnh nhân 13, giá trị
$\alpha=0,5$ có thể tình cờ gần với giá trị tối ưu của bài báo (FGL khớp sát); ở bệnh nhân 3 và
5, giá trị tối ưu thực sự theo bài báo nhiều khả năng khác xa 0,5, nên FGL tái hiện với $\alpha$
cố định không đạt được mức cải thiện tương ứng. Đây là một giả thuyết hợp lý về nguyên nhân
chính của khoảng cách với bài báo, nhưng \textbf{chưa được kiểm chứng trực tiếp} (cần quét
$\alpha$ theo từng bệnh nhân như bài báo gốc để xác nhận).

\section{Các vấn đề phát hiện khi chạy mã nguồn}

Quá trình tái hiện phát hiện chín vấn đề cụ thể trong mã nguồn công khai của tác giả, cho thấy
mã nguồn ở trạng thái công bố \textbf{không chạy được ngay} mà cần chỉnh sửa:

\begin{enumerate}
	\item \textbf{Bốn tệp còn dấu xung đột merge chưa giải quyết} (\texttt{base\_exp.py},
	\texttt{drift\_exp.py}, \texttt{utils/utils.py}, \texttt{utils/generate.py}) - gây
	SyntaxError ngay khi import, mã nguồn không chạy được ở trạng thái công bố.
	\item \textbf{Lệch tầm dự báo trong \texttt{generate.py}}: một trong hai nhánh xung đột làm
	lệch \texttt{forecasting\_horizon} đi 1 bước; đối chiếu với \texttt{data.pkl} tác giả commit
	sẵn xác nhận nhánh đúng.
	\item \textbf{Tham số dòng lệnh của \texttt{drift\_exp.py} bị vô hiệu}: khối
	\texttt{argparse} nằm trong chuỗi bị comment; mã thực sự chạy là đoạn hardcode quét horizon
	2-15.
	\item \textbf{Thiếu khai báo dependency}: \texttt{utils.py} import \texttt{jitcdde} nhưng
	không có trong \texttt{requirements.txt}.
	\item \textbf{Đường dẫn chạy khác tài liệu hướng dẫn}: phải chạy
	\texttt{PYTHONPATH=. python exp/base\_exp.py}, không phải \texttt{python base\_exp.py} như
	tài liệu ghi.
	\item \textbf{Hardcode \texttt{.to("cuda")}} ở ba script CHB-MIT, không chạy được nếu không
	có GPU NVIDIA; đã vá thành tự nhận biết MPS/CUDA/CPU.
	\item \textbf{Nghi vấn \texttt{shuffle=False}} \label{item:shuffle-false} trên dữ liệu
	teacher CHB-MIT đã sắp theo thứ tự lớp (toàn interictal rồi toàn ictal) - có thể là
	nguyên nhân AUC bất thường (0,28) ở lần chạy local, nhưng chưa xác nhận trực tiếp.
	\item \textbf{Cache gom toàn bộ dữ liệu một loại nhãn vào bộ nhớ} trước khi ghi đĩa -
	nguyên nhân hết bộ nhớ khi xử lý Giai đoạn 2/3 của CHB-MIT trên máy cá nhân.
	\item \textbf{Cache đặc trưng STFT (\texttt{.hickle}) không được dọn qua các lần chạy và
	không có cơ chế phục hồi khi ghi dở dang}: trên GPU Kaggle (dung lượng đĩa giới hạn), cache
	của nhiều bệnh nhân cộng dồn làm hết đĩa; khi đó tệp \texttt{.hickle} đang ghi bị cắt giữa
	chừng nhưng vẫn tồn tại trên đĩa, và lần chạy sau cố \textit{đọc lại} tệp hỏng đó gây
	\texttt{UnpicklingError} không được \texttt{try/except} bắt, làm sập toàn bộ tiến trình cho
	bệnh nhân đó thay vì tự tính lại từ đầu. Báo cáo xử lý vấn đề này ở tầng điều phối (notebook
	Kaggle): tự xoá cache của bệnh nhân sau mỗi giai đoạn dùng xong, và thử lại một lần kèm xoá
	cache nếu lần chạy đầu thất bại.
\end{enumerate}

\section{Nhận xét về khả năng tái hiện}

Nhìn chung, thực nghiệm Mackey-Glass cho thấy FGL là một phương pháp \textbf{tái hiện được}
cả về hướng và độ lớn cải thiện, với điều kiện phải sửa các lỗi mã nguồn ở mục trên trước; kết
quả ổn định trên nhiều seed, không phụ thuộc vào một lần chạy may rủi. Trên CHB-MIT, FGL cải
thiện trên ba bệnh nhân có đủ dữ liệu so sánh (1, 3, 5), nhưng mức cải thiện trung bình đo được
lệch đáng kể giữa hai lần chạy độc lập cùng ba bệnh nhân này: $+31,2\%$ ở lần 1 và $+65,2\%$ ở
lần 2 (mục~\ref{sec:chbmit-hailanchay}) - cả hai đều thấp hơn mức 44,8\% bài báo công bố,
nhưng khoảng cách \textit{giữa hai lần chạy của chính báo cáo này} còn lớn hơn khoảng cách tới
con số của bài báo, cho thấy với số bệnh nhân và số trial hiện có, mức cải thiện đo được chưa
đủ ổn định để khẳng định một con số cụ thể. Điểm cải thiện nhất quán cả hai lần là \textit{hướng}
của FGL (luôn tốt hơn hoặc xấp xỉ Baseline, không đảo chiều mạnh) và việc bệnh nhân có Baseline
yếu luôn được FGL cải thiện nhiều hơn bệnh nhân có Baseline đã mạnh. Bệnh nhân 13 (AUC cao nhất
trong bài báo gốc) cho Student/FGL rất cao và ổn định (0,98) ngay ở Giai đoạn 1 và 3, cùng hướng
với bài báo dù chưa có Baseline để so sánh mức tăng. Bệnh nhân 9 vẫn chưa huấn luyện thành công
ở bất kỳ giai đoạn nào sau nhiều lần thử, và bệnh nhân 23 mới có Teacher - phần mở rộng lên
sáu bệnh nhân vì vậy chưa hoàn tất, nên các nhận định trên CHB-MIT trong báo cáo này vẫn nên
được đọc là bằng chứng ban đầu. Điểm tích cực là chỉ số AUC bất thường (0,28) của Teacher trên
máy cá nhân không lặp lại trên Kaggle ở bất kỳ bệnh nhân hay lần chạy nào. Việc không tái hiện
được AES hoàn toàn là do giới hạn thời gian và việc tải dữ liệu, không phản ánh đánh giá về bản
thân phương pháp trên bộ dữ liệu này.
